\section{Viable subspaces in coordinates}

When $\mathcal H_A=\mathbb C^\sigma$ is the space of configurations
$\sigma$ of a finite set of sites and $\mathcal H_B=\mathbb C^\beta$, the
viability of a subspace $S\subseteq\mathcal H_A$ for a target
$T\subseteq\mathcal H_A\otimes\mathcal H_B$~\cite[lines
    433--440]{Arad2016RigorousRG} can be stated in $\mathbb C^{\sigma\times\beta}$,
and the cuts inside $\mathcal H_A$ used by the trimming of a viable
subspace~\cite[lines 660--697]{Arad2016RigorousRG} keep $\mathcal H_B$ on the
right.

\begin{theorem}[Cut coordinates with an environment factor]
    \label{thm:cut_environment_coordinates}
    \lean{Entropy.cutEnvEquiv, Entropy.cutEnvVector,
        Entropy.cutEnvVector_leftApply_of_cutOperator_eq_kronecker_one,
        Entropy.cutEnvVector_leftApply_localLift,
        Entropy.cutEnvVector_leftApply_of_cutOperator_eq_one_kronecker,
        Entropy.IsSupportedOn.schmidtRank_cutEnvVector_leftApply_le,
        Entropy.cutEnvVector_rightApply,
        Entropy.schmidtRank_cutEnvVector_rightApply_le,
        Entropy.schmidtRank_cutVector_slice_le}
    \leanok
    \uses{lem:supported_product_decomposition, def:schmidt_rank}
    For a region $D$ of the sites, write a vector of
    $\mathbb C^\sigma\otimes\mathbb C^\beta$ in the bipartite coordinates
    $\mathbb C^{\sigma_D}\otimes(\mathbb C^{\sigma_{D^{\mathrm c}}}\otimes
        \mathbb C^\beta)$; this is an isometry. An operator $K\otimes I$ on the
    sites acts on the left factor, an operator supported outside $D$ or
    acting on $\mathbb C^\beta$ acts on the right factor and does not
    increase the Schmidt rank across $D$, and a slice $w(\cdot,b)$ has
    Schmidt rank across $D$ at most that of $w$.
\end{theorem}

\begin{proof}
    \leanok
    Each statement is a computation of the coefficient matrix in the cut
    coordinates, using the product form of an operator supported in or
    outside $D$; an operator on the right factor multiplies the coefficient
    matrix on the right, which does not increase its rank.
\end{proof}

\begin{theorem}[Viability in coordinates]
    \label{thm:viable_coordinates}
    \lean{Submodule.isCloseTo_map_iff, Submodule.prodExt,
        Submodule.euclideanTensorEquiv, Submodule.map_prodExt_euclideanTensorEquiv,
        Submodule.isCloseTo_prodExt_iff_isViable}
    \leanok
    \uses{def:subspace_close, thm:subspace_close_iff, def:subspace_viable}
    Closeness of subspaces is invariant under linear isometries. Let
    $S\otimes\mathbb C^\beta\subseteq\mathbb C^{\sigma\times\beta}$ be the
    vectors whose slices $w(\cdot,b)$ lie in $S$. The isometry
    $\mathbb C^{\sigma\times\beta}\cong\mathbb C^\sigma\otimes\mathbb C^\beta$
    sending the coordinate vector of $(a,b)$ to $e_a\otimes e_b$ carries
    $S\otimes\mathbb C^\beta$ onto $S\otimes\mathbb C^\beta$ in the tensor
    product, so $S\otimes\mathbb C^\beta$ is $\delta$-close to $T$ exactly
    when $S$ is $\delta$-viable for the image of $T$.
\end{theorem}

\begin{proof}
    \leanok
    Under an isometry $\Phi$, the projection onto $\Phi U$ is $\Phi P_U\Phi^{-1}$,
    so the vector form of closeness transfers. A vector $w$ with slices in $S$
    is $\sum_b w(\cdot,b)\otimes e_b$, and $x\otimes y$ with $x\in S$ has slices
    $y_b\,x\in S$.
\end{proof}

\begin{theorem}[Matrix units applied to a nonzero subspace]
    \label{thm:basis_matrix_units_generate}
    \lean{Submodule.iSup_map_basis_matrixUnits_eq_top}
    \leanok
    Let $E$ be a vector space over a field $k$, with basis $(b_i)$
    and coordinate functionals $b_j^*$. For every nonzero subspace $S$,
    \begin{align}
        \sum_{i,j}(b_i b_j^*)S=E.
    \end{align}
    The basis need not be finite; the sum denotes the algebraic sum of subspaces.
\end{theorem}
\begin{proof}\leanok
    Choose $x\in S\setminus\{0\}$ and $j$ with $b_j^*(x)\ne0$.
    Each basis vector $b_i$ is a scalar multiple of $(b_i b_j^*)x$.
    Thus the sum contains a basis and equals $E$.
\end{proof}

\begin{theorem}[A viable subspace for a nonzero target is nonzero]
    \label{thm:viable_nonzero_target}
    \lean{Submodule.isViable_ne_bot_of_target_ne_bot}
    \leanok
    \uses{def:subspace_viable}
    Let $E,F$ be finite-dimensional complex inner-product spaces.
    If $T\subseteq E\otimes F$ is nonzero and $S$ is $\delta$-viable for $T$
    with $\delta<1$, then $S$ is nonzero.
\end{theorem}
\begin{proof}\leanok
    \uses{def:subspace_viable,thm:subspace_close_iff}
    If $S=0$, the projection onto $S\otimes F$ vanishes. Its required lower
    bound on a nonzero vector of $T$ contradicts $1-\delta>0$.
\end{proof}

\begin{theorem}[Extending a finite operator family by zero]
    \label{thm:operator_family_zero_extension}
    \lean{Submodule.exists_clm_family_extension_of_card_le}
    \leanok
    Let $E$ be a complex normed space and $(A_i)_{i\in I}$ a finite family
    of continuous linear maps on $E$. If $|I|\le M$, there are an injection
    $f:I\to\{0,\ldots,M-1\}$ and maps $B_j$ with $B_{f(i)}=A_i$ and
    $B_j=0$ outside the image of $f$. For every subspace $S\subseteq E$,
    \begin{align}
        \sum_{j=0}^{M-1}B_jS=\sum_{i\in I}A_iS.
    \end{align}
    The family $B$ is chosen before $S$. Neither completeness nor finite
    dimension of $E$ is required.
\end{theorem}
\begin{proof}\leanok
    Enumerate $I$ inside the first $M$ integers, place $A_i$ at its image,
    and put zero elsewhere. Every image subspace on either side is then
    either an original image or the zero subspace.
\end{proof}

\begin{lemma}[Normalized projections from close subspaces]
    \label{lem:close_normalized_projection}
    \lean{Submodule.IsCloseTo.exists_unit_mem_norm_inner_sq_ge}
    \leanok
    \uses{thm:subspace_close_iff}
    Let $S,T$ be subspaces of a finite-dimensional complex Hilbert space.
    Suppose $S$ is $\delta$-close to $T$, where $\delta<1$.
    For every unit vector $\psi\in T$, put $p=P_S\psi$.
    Then $p\ne0$, and $\phi=p/\|p\|$ belongs to $S$, has norm one, and satisfies
    \begin{align}
        \langle\phi,\psi\rangle&=\|p\|, &
        |\langle\phi,\psi\rangle|^2&\ge1-\delta, &
        \|\phi-\psi\|^2&\le2\delta.
    \end{align}
    In particular, the phase of $\phi$ is aligned with $\psi$.
\end{lemma}
\begin{proof}\leanok
    Closeness gives $1-\delta\le\|p\|^2$, so $p\ne0$.
    Orthogonality gives $\langle p,\psi\rangle=\|p\|^2$.
    Since $\|p\|\le1$, one has $\|p\|^2\le\|p\|$ and hence
    $\|\phi-\psi\|^2=2-2\|p\|\le2\delta$.
\end{proof}

\begin{theorem}[Unit approximants of bounded Schmidt rank]
    \label{thm:viable_unit_schmidt_approximation}
    \lean{Submodule.IsViable.exists_unit_mem_prodExt_schmidtRank_le}
    \lean{Submodule.IsViable.exists_unit_schmidtRank_le_of_equiv}
    \leanok
    \uses{def:subspace_viable, def:schmidt_rank, thm:viable_coordinates}
    Let $\sigma,\beta$ be finite sets, and let
    $\mathcal U:\C^{\sigma\times\beta}\to\C^\sigma\otimes\C^\beta$
    be the isometry sending $e_{(a,b)}$ to $e_a\otimes e_b$.
    Suppose $S\subseteq\C^\sigma$ is $\delta$-viable for
    $T\subseteq\C^\sigma\otimes\C^\beta$, where $\delta<1$.
    Every unit vector $\psi\in\C^{\sigma\times\beta}$ with
    $\mathcal U\psi\in T$ admits a unit vector
    $\phi\in\C^{\sigma\times\beta}$ such that every column
    $\phi(\cdot,b)$ lies in $S$ and
    \begin{align}
        |\langle\phi,\psi\rangle|^2 & \ge1-\delta, \\
        \SR(\phi)                   & \le\dim S.
        \label{eq:viable_unit_schmidt_bounds}
    \end{align}
    An equivalent formulation uses a complex inner-product space $E$
    and a specified linear isometric isomorphism
    $e:E\to\C^{\sigma\times\beta}$.
    If $T_0\subseteq E$, $\delta<1$ and $S$ is $\delta$-viable for
    $(\mathcal U\circ e)T_0$, then every unit vector $\psi_0\in T_0$
    admits a unit vector $\phi_0\in E$ such that
    $(e\phi_0)(\cdot,b)\in S$ for every $b\in\beta$,
    $|\langle\phi_0,\psi_0\rangle|^2\ge1-\delta$ and
    $\SR(e\phi_0)\le\dim S$.
    The rank in this formulation is measured in the product coordinates
    specified by $e$.
    These are the unit-vector consequences of viability used
    in~\cite[lines 1080--1094]{Arad2016RigorousRG}.
\end{theorem}

\begin{proof}
    \leanok
    \uses{thm:viable_coordinates, lem:close_normalized_projection}
    Let $\mathcal S$ be the subspace of arrays whose columns lie in $S$,
    and put $p=P_{\mathcal S}\psi$. Viability in coordinates gives
    $\|p\|^2\ge1-\delta>0$. The normalized projection
    $\phi=p/\|p\|$ belongs to $\mathcal S$ and satisfies
    $|\langle\phi,\psi\rangle|^2=\|p\|^2$.
    The column space of its coefficient matrix is contained in $S$,
    so its rank is at most $\dim S$.
    For the formulation on $E$, apply the coordinate assertion to
    $e\psi_0$, obtaining $\chi$, and set $\phi_0=e^{-1}\chi$.
    Then $\|\phi_0\|=\|\chi\|=1$ and
    $\langle\phi_0,\psi_0\rangle=\langle\chi,e\psi_0\rangle$.
    The column and rank assertions follow from $e\phi_0=\chi$.
\end{proof}

\begin{theorem}[Leading marginal mass from viability]
    \label{thm:viable_leading_marginal_mass}
    \lean{Submodule.IsViable.one_sub_le_kyFanNorm}
    \leanok
    \uses{def:subspace_viable, def:ky_fan_norm, thm:viable_coordinates}
    In the coordinate formulation of
    Theorem~\ref{thm:viable_unit_schmidt_approximation}, let
    $R=\dim S$, let $C_\psi(a,b)=\psi(a,b)$, and set
    $\rho_\psi=C_\psi C_\psi^\dagger$.
    The sum of the largest $R$ eigenvalues of $\rho_\psi$ satisfies
    $\|\rho_\psi\|_{(R)}\ge1-\delta$.
\end{theorem}

\begin{proof}
    \leanok
    \uses{thm:viable_unit_schmidt_approximation,
        thm:homogeneous_schmidt_overlap_bound}
    Choose the unit vector $\phi$ from the coordinate formulation of
    Theorem~\ref{thm:viable_unit_schmidt_approximation}.
    The maximal-overlap bound gives
    $1-\delta\le|\langle\psi,\phi\rangle|^2
        \le\|\rho_\psi\|_{(R)}$.
\end{proof}

\begin{theorem}[Tensor exchange in product coordinates]
    \label{thm:tensor_exchange_product_coordinates}
    \lean{Submodule.euclideanTensorEquiv_symm_comm}
    \leanok
    \uses{thm:viable_coordinates}
    Let $\sigma,\beta$ be finite sets. Write
    $E_{\sigma,\beta}:\C^{\sigma\times\beta}\to
    \C^\sigma\otimes\C^\beta$ for the isometry sending
    $e_{(a,b)}$ to $e_a\otimes e_b$, and let
    $\Sigma(x\otimes y)=y\otimes x$.
    For every $z\in\C^{\beta\times\sigma}$,
    \begin{align}
        \bigl(E_{\sigma,\beta}^{-1}\Sigma E_{\beta,\sigma}z\bigr)(a,b)
        &=z(b,a).
        \label{eq:tensor_exchange_coordinates}
    \end{align}
\end{theorem}
\begin{proof}
    \leanok
    Evaluate both sides on $z=e_{(b,a)}$ and extend linearly.
\end{proof}

\section{Canonical low-energy spaces under scalar shifts}

The following identities identify the original low-energy spaces used in
the low-density construction of \cite{Arad2016RigorousRG}. All scalar
normalizations of one Hermitian matrix lead to the same original windows.

\begin{lemma}[Independence of the complete eigenbasis]
    \label{lem:closed_spectral_basis_independence}
    \lean{Submodule.spectralSpanLE_eq_of_eigenbases}
    \leanok
    Let $H$ be a linear operator on a real or complex inner-product space.
    Suppose $b$ and $b'$ are complete finite orthonormal eigenbases of $H$,
    with real eigenvalues $a_i$ and $a'_j$. For every real $c$,
    \begin{align*}
      \operatorname{span}\{b_i:a_i\le c\}
       =\operatorname{span}\{b'_j:a'_j\le c\}.
    \end{align*}
    Repeated eigenvalues and the zero-dimensional space are included.
\end{lemma}
\begin{proof}
    \leanok
    Expand $b_i$ in the second basis. A coefficient at $b'_j$ is zero
    whenever $a_i\ne a'_j$, by applying $H$ to the expansion. Thus an
    eigenvector with $a_i\le c$ has zero coefficients at every $a'_j>c$.
    The reverse inclusion follows by exchanging the bases.
\end{proof}

\begin{lemma}[A scalar shift of a spectral cutoff]
    \label{lem:closed_spectral_scalar_shift}
    \lean{Submodule.spectralSpanLE_sub}
    \leanok
    For a complete finite orthonormal basis $b$, real numbers $a_i$, and
    real $s,c$, one has
    \begin{align*}
      \operatorname{span}\{b_i:a_i-s\le c-s\}
        =\operatorname{span}\{b_i:a_i\le c\}.
    \end{align*}
\end{lemma}
\begin{proof}
    \leanok
    The two inequalities select precisely the same indices.
\end{proof}

\begin{theorem}[The canonical spectral space after a scalar shift]
    \label{thm:canonical_spectral_scalar_shift}
    \lean{Matrix.IsHermitian.spectralSpanLE_sub_smul_one}
    \leanok
    Let $H$ be a finite Hermitian matrix and $s,c\in\mathbb R$.
    The canonical spectral space of $H-s\Id$ below $c-s$ equals the
    canonical spectral space of $H$ below $c$. No relation between
    independently selected eigenbases is assumed.
\end{theorem}
\begin{proof}
    \leanok
    \uses{lem:closed_spectral_basis_independence,lem:closed_spectral_scalar_shift}
    Every eigenvector of $H$ with eigenvalue $a_i$ is an eigenvector of
    $H-s\Id$ with eigenvalue $a_i-s$. Apply independence of the complete
    eigenbasis, followed by the scalar-cutoff identity.
\end{proof}

\begin{corollary}[Original windows from actual centered ground data]
    \label{cor:canonical_centered_ground_window}
    \lean{Matrix.IsHermitian.spectralSpanLE_centered_ground}
    \lean{Matrix.IsHermitian.spectralSpanLE_centered_ground_of_eq}
    \leanok
    Let $H$ be a Hermitian matrix on a nonzero finite-dimensional space.
    Suppose $\|z\|=1$, $(H-s\Id)z=E_0z$ and
    $H-(s+E_0)\Id\ge0$. Then
    $\min\operatorname{spec}H=E_0+s$. For every $\eta\in\mathbb R$,
    the canonical spectral space of $H-s\Id$ below $E_0+\eta$ equals
    the canonical spectral space of $H$ below
    $\min\operatorname{spec}H+\eta$. The same conclusion holds for
    the canonical spectral space of any Hermitian presentation $B=H-s\Id$.
\end{corollary}
\begin{proof}
    \leanok
    \uses{thm:canonical_spectral_scalar_shift}
    The normalized eigenvector attains the lower bound $E_0+s$, so that
    number is the original minimum eigenvalue. Apply the preceding theorem
    with cutoff $\min\operatorname{spec}H+\eta$.
\end{proof}


\begin{lemma}[A regional cut with the full exterior]
    \label{thm:regional_cut_full_exterior}
    \lean{Entropy.subregionCut, Entropy.mem_subregionCut,
        Entropy.subregionCutSubtypeEquiv, Entropy.subregionRestSubtypeEquiv,
        Entropy.regionConfigurationSplit,
        Entropy.schmidtRank_cutEnvVector_regionConfigurationSplit}
    \leanok
    \uses{thm:cut_environment_coordinates}
    Let $A\subseteq J\subseteq V$ be finite sets of sites with common
    on-site dimension $q$. Write $F_J\psi$ for the configuration vector
    of $\psi$ in $\mathcal H_J\otimes\mathcal H_{J^c}$. Splitting the
    first factor at $A$ and retaining the entire exterior gives
    \begin{align}
        \SR_{A:(J\setminus A)\cup J^c}(F_J\psi)
        =\operatorname{rank}\rho_A(\psi).
        \label{eq:regional_cut_full_exterior}
    \end{align}
    No exterior factor is traced out before computing the Schmidt rank.
\end{lemma}
\begin{proof}\leanok
    The two coefficient matrices differ only by the row identification of
    $A$ and the column identification
    $(J\setminus A)\sqcup J^c\cong A^c$. Reindexing preserves rank, and
    the rank of a coefficient matrix equals that of its product with its
    conjugate transpose.
\end{proof}


\begin{lemma}[Actual regional tensor coordinates]
    \label{thm:regional_tensor_viability}
    \lean{Entropy.regionTensorCoordinates,
        Entropy.isCloseTo_prodExt_region_iff_isViable,
        Entropy.regionTensorCoordinates_eq_exteriorFirst}
    \leanok
    \uses{thm:viable_coordinates,thm:regional_cut_full_exterior}
    Let $G\subseteq\mathcal H_V$ and $S\subseteq\mathcal H_J$.
    The coordinate extension $S\otimes\mathcal H_{J^c}$ is
    $\delta$-close to $F_JG$ if and only if $S$ is $\delta$-viable for
    the image of $G$ under the canonical tensor isometry
    $\mathcal H_V\cong\mathcal H_J\otimes\mathcal H_{J^c}$.
    Splitting the exterior first and then exchanging the tensor factors
    gives the same isometry.
\end{lemma}
\begin{proof}\leanok
    The Euclidean tensor isometry carries a coordinate extension onto the
    tensor product subspace. Both orders of the configuration split evaluate
    $\psi$ at the same global configuration, so they agree after exchanging
    the tensor factors.
\end{proof}


\begin{lemma}[Tensor factor changes for viable spaces]
    \label{thm:viable_tensor_factor_changes}
    \lean{Submodule.map_tensorProd_congrIsometry,
        Submodule.map_tensorProd_commIsometry, Submodule.tensorProd_top_top,
        Submodule.map_tensorProd_assocIsometry, Submodule.finrank_tensorProd_le,
        Submodule.isViable_map_congrIsometry_iff,
        Submodule.isViable_map_assocIsometry_iff,
        Submodule.IsViable.tensorProd_of_factor_splits,
        Submodule.euclideanTensorEquiv_reindex_symm_tmul}
    \leanok
    \uses{def:subspace_viable,thm:viable_tensoring,thm:viable_coordinates}
    Let $E,F,R$ be finite-dimensional real or complex inner-product spaces.
    Tensor product subspaces are carried to tensor product subspaces by
    isometries of their factors, by exchange of factors, and by reassociation.
    Their dimension is at most the product of the input dimensions.
    Viability is invariant under isometries of the retained factor and the
    environment. For one fixed target $T\subseteq(E\otimes F)\otimes R$,
    if $S_E$ is $\delta_E$-viable in its first-factor split and $S_F$ is
    $\delta_F$-viable in its second-factor split, then $S_E\otimes S_F$
    is $(\delta_E+\delta_F)$-viable for $T$.
    For complex coordinate spaces and a finite reindexing $e:\kappa\cong\sigma\times\beta$,
    the inverse tensor isometry sends $x\otimes y$ to
    $c\mapsto x(e(c)_1)y(e(c)_2)$.
\end{lemma}
\begin{proof}\leanok
    Tensor product subspaces are generated by their pure tensors.
    Transport each generator and extend linearly. The dimension bound
    follows from the range of the tensor product of the two inclusions.
    Transport the two viability inequalities to the same three-factor
    space and apply tensoring there. The coordinate formula is the
    product formula for pure tensors after the specified reindexing.
\end{proof}


\begin{lemma}[The two regional tensor diagrams]
    \label{thm:regional_tensor_diagrams}
    \lean{Entropy.unionTensorCoordinates, Entropy.exteriorTensorCoordinates,
        Entropy.rightExteriorTensorCoordinates,
        Entropy.rightExteriorTensorCoordinates_symm_tmul_apply,
        Entropy.regionTensorCoordinates_symm_tmul_apply,
        Entropy.unionTensorCoordinates_symm_tmul_apply,
        Entropy.exteriorTensorCoordinates_symm_tmul_apply,
        Entropy.threeRegionTensorCoordinates,
        Entropy.threeRegionTensorCoordinates_assoc,
        Entropy.threeRegionTensorCoordinates_comm_assoc}
    \leanok
    \uses{thm:regional_tensor_viability,thm:viable_tensor_factor_changes}
    Let $A,B\subseteq V$ be disjoint and write $R=(A\cup B)^c$.
    The actual configuration restrictions define isometries
    $\mathcal U:\mathcal H_{A\cup B}\to\mathcal H_A\otimes\mathcal H_B$,
    $\mathcal H_{A^c}\to\mathcal H_B\otimes\mathcal H_R$ and
    $\mathcal H_{B^c}\to\mathcal H_A\otimes\mathcal H_R$.
    Let $\mathcal C:\mathcal H_V\to
        (\mathcal H_A\otimes\mathcal H_B)\otimes\mathcal H_R$
    first split the union from $R$ and then apply $\mathcal U$.
    Reassociating $\mathcal C$ equals the actual split at $A$ followed by
    the indicated split of $A^c$. Exchanging $A,B$ and reassociating equals
    the actual split at $B$ followed by the indicated split of $B^c$.
\end{lemma}
\begin{proof}\leanok
    Compare the inverse maps on $x\otimes(y\otimes z)$.
    In each case the value at a global configuration $c$ is the product
    $x(c|_A)y(c|_B)z(c|_R)$, in the appropriate order. Associativity and
    commutativity of scalars give equality. Extend linearly from pure
    tensors; the whole exterior $R$ is retained.
\end{proof}


\begin{lemma}[Merging the actual regional viable spaces]
    \label{thm:actual_regional_viability_merge}
    \lean{Entropy.regionTensorMerge, Entropy.isViable_regionTensorMerge,
        Entropy.finrank_regionTensorMerge_le, Entropy.isViable_region_congr}
    \leanok
    \uses{thm:regional_tensor_diagrams,thm:viable_tensor_factor_changes}
    Let $A,B\subseteq V$ be disjoint, and fix one arbitrary global target
    $G\subseteq\mathcal H_V$. Suppose $S_A\subseteq\mathcal H_A$ is
    $\delta_A$-viable for $G$ in the actual $A:A^c$ split and
    $S_B\subseteq\mathcal H_B$ is $\delta_B$-viable in the actual
    $B:B^c$ split. Then
    $S_{A\cup B}=\mathcal U^{-1}(S_A\otimes S_B)$ is
    $(\delta_A+\delta_B)$-viable for the same $G$ in the actual
    $(A\cup B):(A\cup B)^c$ split, and
    $\dim S_{A\cup B}\le(\dim S_A)(\dim S_B)$.
    The target may be entangled across all three factors.
    Equal regions identify these viable spaces by their canonical
    configuration equivalence.
\end{lemma}
\begin{proof}\leanok
    The two proved tensor diagrams carry both hypotheses to the one
    target $\mathcal CG$. Tensor there and transport the actual candidate
    back through $\mathcal U$. The dimension bound is preserved by this
    isometry. For equal regions, the canonical configuration equivalence
    is the identity after identifying the regions.
\end{proof}

\begin{lemma}[A regional operator in its actual tensor coordinates]
    \label{lem:regional_operator_tensor_coordinates}
    \lean{Entropy.localLift_middleFirst_conjugation,
        Entropy.regionTensorCoordinates_conjugation_localLift}
    \leanok
    \uses{thm:regional_tensor_viability,thm:middle_first_conjugated_tensor_sum}
    Let $V$ be a finite set of sites, $q\ge0$ and $J\subseteq V$.
    For any matrix $K$ on $\mathbb C^{J\to\{1,\ldots,q\}}$, let
    $\widehat K$ be its identity extension to the full configuration space.
    The canonical regional tensor isometry $T_J$ satisfies
    \begin{align}
       T_J\widehat K T_J^{-1}=K\otimes I_{J^c}.
    \end{align}
    The explicit coordinate map obtained by placing the exterior first
    and then exchanging the tensor factors gives the same identity.
    Empty regions and zero local dimensions are included; positivity
    is not required.
\end{lemma}
\begin{proof}\leanok
    In the exterior-first configuration split, the entries of
    $\widehat K$ are those of $I_{J^c}\otimes K$: they vanish when
    the two exterior configurations differ, and otherwise equal the
    corresponding entry of $K$. Apply the conjugated tensor-sum identity
    to this singleton sum. Exchange of the two tensor factors gives
    $K\otimes I_{J^c}$, and the two standard coordinate constructions
    are the same isometry.
\end{proof}
